\section{Conclusion, Discussion, and Limitations}
\label{sec:conclusion}

\subsection{Summary}

This project gave empirical content to the self-resolving prediction market of
\textcite{srinivasan2025selfresolving}, a mechanism that elicits and aggregates beliefs
about an \emph{unverifiable} outcome by paying agents the cross-entropy of their report
against a late reference agent rather than against the truth. We reproduced the
mechanism in an open-source agent-based simulator and evaluated it along two axes. In
the clean synthetic world (RQ1), a freely deviating agent's profit-maximising
report converges onto truthful reporting as the reference accumulates informational
substitutes ($\gamma^\star \to 1$ by $k \approx 64$), operationalising the paper's
Theorems~1--3; we also reproduced the analytic minimum-substitutes bound $k_{\min}$ and
confirmed exact aggregation under truthful play. On real Polymarket data (RQ2), the
self-resolving payouts track the verifiable payouts they replace almost perfectly at a
genuinely late reference (Pearson $\approx 1.00$, payout ratio $\approx 0.95$), and
degrade gracefully as the reference moves earlier, isolating outcome leakage as the
channel driving the residual gap.

\subsection{Discussion}

The two studies answer two genuinely different questions, and it is worth being precise
about which is which. RQ1 is about \emph{incentives}: in a world where the paper's
assumptions hold, a rational agent who is free to lie chooses not to, once the reference
is informed enough. RQ2 is about \emph{outputs}: on real data where the assumptions fail,
the payouts the mechanism produces still resemble the payouts an outcome-based mechanism
would have produced. Neither subsumes the other --- RQ2 is not an incentive test, and RQ1
is not a test on real behaviour --- but together they bracket the mechanism from both
sides and suggest it is more robust in practice than its worst-case theory promises (the
empirical onset of truthfulness at $k\approx 64$ is well inside the $k_{\min}\approx 146$
worst-case bound).

The most interesting tension the empirical study surfaces is that the mechanism's
residual error lives almost entirely in the \emph{ranking} of who deserves the most
reward, not in the payout \emph{levels}. A self-resolving designer can recover roughly
the right total spend and reward roughly the right behaviour at an early, forecast-only
reference; what they cannot fully recover without leaking the outcome is the fine-grained
ordering of which trader moved the price toward truth. This points directly at a design
question: can the mechanism be adjusted so that the ranking better matches who actually
moved the price toward truth, without needing the outcome?

\subsection{Limitations}

Several limitations qualify these results and frame future work.
\begin{itemize}
    \item \textbf{RQ2 measures outputs, not incentives.} We replay fixed historical
    price paths and never let agents deviate; the empirical study therefore says nothing
    about whether real traders would report truthfully under the live mechanism. A live
    or agent-in-the-loop deployment would be needed for that.
    \item \textbf{Broken assumptions on real data.} Real markets violate essentially all
    of A1--A5: traders are not risk-neutral Bayesians, there is no common prior, and price
    moves are serially correlated (herding, momentum, shared news) rather than
    conditionally independent given $Y$. The agreement we observe is thus evidence of
    \emph{robustness}, not of the assumptions holding.
    \item \textbf{The reference conflates two things.} A late reference price is both
    better-informed and outcome-contaminated; our sweeps separate these channels only
    indirectly, by varying lateness. A cleaner instrument for ``well-informed but not
    omniscient'' would strengthen the claim.
    \item \textbf{Scope of the synthetic model.} The simulator deliberately makes A1--A4
    hold by construction and only varies A5. This is the right design for isolating the
    core incentive claim, but it cannot speak to how the mechanism behaves when, say,
    conditional independence is mildly violated --- a natural next experiment.
    \item \textbf{Binary outcomes only.} Both the mechanism as implemented and our
    studies are restricted to binary $Y$. Multi-outcome and continuous predictions are
    left for future work.
\end{itemize}

\subsection{Future work}

Two directions follow directly from the limitations. First, a \emph{mechanism-design}
question: can the CE-MSR be reweighted or the reference reconstructed so that the
\emph{ranking} of rewards tracks who pushed the price toward truth, closing the gap that
RQ2 attributes to leakage --- without observing the outcome? Second, a \emph{generality}
question: extending the simulator and the empirical replay to non-binary outcomes, and
to information structures that depart from strict conditional independence, would test how
far the reference-as-proxy idea stretches beyond the regime where it is provably
truthful. More broadly, the gap between the loose worst-case $k_{\min}$ and the much
earlier empirical onset of truthfulness suggests there is room for tighter,
prior-dependent guarantees that better match observed behaviour.
